\documentclass[11pt]{article}
\usepackage[T1]{fontenc}
\usepackage{lmodern}
\usepackage{amsmath,amssymb,amsthm}
\usepackage[margin=1.1in]{geometry}
\usepackage{microtype}
\usepackage[hidelinks]{hyperref}

\newtheorem{theorem}{Theorem}

\title{Failure of automatic mass conservation for\\
renormalized nonlinear reaction--diffusion}
\author{Kaiwen Jin\thanks{\href{mailto:yui.ui.siki@gmail.com}{\texttt{yui.ui.siki@gmail.com}}}}
\date{27 November 2026}
\hypersetup{pdfauthor={Kaiwen Jin},pdftitle={Failure of automatic mass conservation for renormalized nonlinear reaction-diffusion}}

\begin{document}
\maketitle

\begin{abstract}
The truncated formulation of a nonlinear reaction--diffusion system does
not by itself imply conservation of every reaction invariant. We construct
a five-component system on an interval with a locally Lipschitz,
quasi-positive, entropy-controlled reaction. The solution satisfies the
full truncated identity and the required diffusion energy bounds, while
the mass associated with a reaction invariant exceeds its prescribed
initial value. The excess mass concentrates at the boundary as
\(t\downarrow0\), while the bounded renormalizations have the prescribed
initial trace. This
gives a counterexample to AIM Problem~483.
\end{abstract}

\section{Introduction}

For a classical reaction--diffusion solution with no flux through the
boundary, a linear combination of the species masses is conserved whenever
the corresponding linear combination of the reaction rates vanishes.
For renormalized solutions, the admissible tests have compactly supported
derivatives in the concentration variables. The linear mass test is outside
this class, and the cutoff limit must be justified. For linear diffusion,
Fischer's author preprint gives conservation in the no-flux case
under its entropy-dissipation and solution hypotheses
\cite[Proposition~6 in arXiv:1703.00730v1]{Fischer}.

Fellner, Fischer, Kniely and Tang construct global renormalized solutions
with nonlinear diffusion \cite[Theorem~1.3]{FFKT}. Their solution definition
includes the truncated identity, an entropy inequality and conservation
laws. They ask whether the truncated identity alone implies the last two
properties \cite[Remark~1.2]{FFKT}. AIM Problem~483 poses the question for
the conservation laws with all diffusion exponents in \((0,2)\)
\cite{AIM483}.

We give a negative answer. In the example, a concentrating first species
has finite diffusion energy. A second species with a broader heat profile
supplies the diffusion potential required by the balance for their sum.
The excess mass in their sum concentrates at \(x=0\) as \(t\downarrow0\).
The prescribed datum has finite entropy, and tests with compactly
supported concentration derivatives do not detect the excess.
A third species permits the pointwise entropy bound on the reaction. Two bounded components make
the prescribed reaction residuals functions of the concentrations.
The construction is explicit near the initial time, and the existence
theorem of \cite{FFKT} supplies a global continuation.

\section{The truncated formulation}

Let \(\Omega\subset\mathbb R^d\) be smooth and bounded. The system is
\begin{equation}
 \partial_t u_i=d_i\Delta u_i^{m_i}+f_i(u),\qquad
 \partial_\nu u_i^{m_i}=0,\qquad 1\le i\le N ,
 \label{eq:system}
\end{equation}
where \(N\ge 2\), \(d_i>0\), and \(0<m_i<2\), with at least one \(m_i\ne1\).
The map \(f:[0,\infty)^N\to\mathbb R^N\) is locally Lipschitz and \(f_i(z)\ge0\)
when \(z_i=0\). We require
\begin{equation}
 \sum_i f_i(z)(\log z_i+\mu_i)
 \le C\sum_i(1+z_i\log z_i),\qquad z\in(0,\infty)^N ,
 \label{eq:entropy}
\end{equation}
for fixed \(\mu_i\in\mathbb R\) and \(C\ge0\). The nonnegative datum \(u_0\)
satisfies \(\sum_i\int_\Omega u_{i,0}(1+|\log u_{i,0}|)<\infty\),
with the logarithm-weighted factor set to zero at zero.

Put \(Q_T=(0,T)\times\Omega\) and \(w_i=u_i^{m_i/2}\). The regularity
requirements are \(u_i\ge0\), \(u_i\in L^\infty(0,T;L^1(\Omega))\), and
\(w_i\in L^2(0,T;H^1(\Omega))\) for every \(T<\infty\).
For every smooth \(\xi\) on \([0,\infty)^N\) with compactly supported \(D\xi\),
every smooth space-time test \(\psi\) on the closed cylinder, and almost every
terminal time \(T>0\), the truncated identity reads
\begin{equation}
\begin{aligned}
 &\int_\Omega \xi(u(T))\psi(T)-\int_\Omega \xi(u_0)\psi(0)
       -\int_{Q_T}\xi(u)\partial_t\psi\\
 &\quad=-\sum_{i,j}\frac{4d_i}{m_j}\int_{Q_T}
       \psi\,\xi_{ij}(u)w_i u_j^{1-m_j/2}\nabla w_i\cdot\nabla w_j\\
 &\qquad\phantom{=}-\sum_i2d_i\int_{Q_T}\xi_i(u)w_i\nabla w_i\cdot\nabla\psi
       +\sum_i\int_{Q_T}\xi_i(u)f_i(u)\psi .
\end{aligned}
\label{eq:renormalized}
\end{equation}
Here \(\xi_i,\xi_{ij}\) are derivatives in the concentration variables.
The flux and mixed gradient terms of \eqref{eq:system} are written
in terms of \(w_i\) in \eqref{eq:renormalized}. Tests have unrestricted
boundary values. The solution requirements consist of these spaces and
\eqref{eq:renormalized}.

\begin{theorem}\label{thm:counterexample}
There are data and a solution of \eqref{eq:renormalized} under the stated
hypotheses with \(\Omega=(0,1)\), \(N=5\), \(d_i=1\), and
\((m_1,m_2,m_3,m_4,m_5)=(1/8,1,1,1,1)\). The vector \(q=(1,1,0,0,0)\)
satisfies \(q\cdot f(z)=0\) on the whole nonnegative cone. For some
\(t_*>0\) and \(M>0\), the solution has
\[
 \int_0^1\bigl(u_1(t,x)+u_2(t,x)\bigr)\,dx
 =\int_0^1\bigl(u_{1,0}(x)+u_{2,0}(x)\bigr)\,dx+M
 \qquad (0<t\le t_*) .
\]
\end{theorem}

\section{Profiles for small positive times}

We write \(m=1/8\). Estimates in this section hold for all sufficiently
small \(t>0\), with constants independent of \(t\) and \(x\). The symbols
\(\lesssim\) and \(\asymp\) denote a bound by a fixed positive multiple and
the corresponding two-sided comparison.

\subsection{A heat background and a pair of concentrations}

Take \(p=3/4\), \(\beta=p/2=3/8\), and \(b_0(x)=1+x^{-p}\) on \((0,1)\).
Let \(B\) solve the Neumann heat equation with this datum. Gaussian
convolution of the even, \(2\)-periodic extension of \(b_0\) gives \(B\)
explicitly. It is smooth for \(t>0\), even at both endpoints, and at least
\(2\). The Neumann kernel has integral one in each variable, so \(B\)
has mass \(\int_0^1 b_0\). At every \(x\in(0,1)\), \(B(t,x)\to b_0(x)\).
The datum has finite entropy since \(p<1\).

For \(0\le x\le\sqrt t\), we have
\begin{equation}
 B(t,x)\asymp t^{-\beta},\qquad
 |B_x(t,x)|\lesssim t^{-\beta-1/2},\qquad
 |B_t(t,x)|\lesssim t^{-\beta-1}.
 \label{eq:heatbounds}
\end{equation}
For the lower bound, integrate the Gaussian against \(z^{-p}\) on
\(0<z<\sqrt t\). In the upper bound, the contribution from \(|z|<1\)
is estimated by \(z=\sqrt t\,s\), with \(x/\sqrt t\) in \([0,1]\).
The other cusps of the periodic extension give an exponentially small
contribution, by the Gaussian bound on successive periods. The \(x\)
and \(t\) derivatives of the Gaussian are bounded by the Gaussian at
twice the time times \(Ct^{-1/2}\) and \(Ct^{-1}\), respectively.
In addition,
\begin{equation}
 \int_0^T\!\int_0^1 |\partial_x\sqrt B|^2\,dx\,dt<\infty .
 \label{eq:heatenergy}
\end{equation}
Indeed, on a positive-time strip the entropy \(\int_0^1 B(\log B-1)\,dx\)
has derivative \(-4\int_0^1|\partial_x\sqrt B|^2\,dx\).
Jensen's inequality for the Neumann kernel bounds the entropy by its
finite initial value. The entropy is bounded below by \(-1\), which
gives \eqref{eq:heatenergy} as the lower end of the strip tends to zero.

We prescribe the sum of the first two concentrations and the sum of their
diffusion potentials. Let \(\alpha=3/4\), \(L=t^\alpha\), and \(y=x/L\).
Choose the smooth even bump
\[
 \phi(y)=
 \begin{cases}
 \exp\!\bigl(-1/(1-y^2)\bigr),& |y|<1,\\
 0,& |y|\ge1 .
 \end{cases}
\]
On \(0\le y\le1\) set \(\Phi(y)=\int_y^1 s\phi(s)\,ds\), set it to zero
for \(y\ge1\), and extend it evenly. Then \(\Phi'=-y\phi\) and
\(0\le\Phi\le\phi/2\) for \(y\ge0\). Define
\begin{equation}
 k=L^{-1}\phi(y),\qquad P=\dot L\,\Phi(y),\qquad
 M=\int_0^1\phi(y)\,dy .
 \label{eq:spike}
\end{equation}
Both profiles are supported on \(0\le x\le L\).
We have \(\int_0^1 k\,dx=M>0\) and \(k_t=P_{xx}\) since both
derivatives equal \(-\dot L L^{-2}(\phi+y\phi')\).

Our first two concentrations \(U,W\) will have sum \(B+1+k\) and
diffusion potentials with sum \(B+1+P\). Their difference determines
\(U-U^m=k-P\). Since \(0<m<1\), the function \(s\mapsto s-s^m\)
is strictly increasing from \(0\) to infinity on \(s\ge1\). Denote
its inverse by \(g\). The inequality \(L\dot L\to0\) gives \(k-P\ge k/2\)
for small \(t\), and we set
\begin{equation}
 U=g(k-P),\qquad W=B+1+P-U^m .
 \label{eq:pair}
\end{equation}
Outside the support of \(k\), \(U=1\) and \(W=B\).
On the support, \(U\asymp 1+k\), \(P=O(t^{\alpha-1})\), and
\(U^m=O(t^{-\alpha m})\). Both \(P\) and \(U^m\) are \(o(t^{-\beta})\),
since \(\alpha+\beta>1\) and \(\beta>\alpha m\).
Consequently \(W\ge B/2\) and \(W\ge2\) on the support for sufficiently
small \(t\).

By \eqref{eq:pair}, \(U+W=B+1+k\) and \(U^m+W=B+1+P\).
The equations \(B_t=B_{xx}\) and \(k_t=P_{xx}\) give
\[
 (U+W)_t=(U^m+W)_{xx}.
\]
Thus the individual residuals are \(F=U_t-(U^m)_{xx}\) and \(-F\).
The function \(F\) is supported on \(0\le x\le L\).
The profiles in \eqref{eq:pair} are smooth for \(t>0\) and even at
both endpoints.

\subsection{The integrated entropy residual}

Let
\[
 \ell=\log B(t,0),\qquad e=F\log(U/W),\qquad E(t)=\int_0^1 e(t,x)\,dx .
\]
These are the entropy contributions from the first two species.
We will use both a pointwise bound on \(e\) and the limit
\begin{equation}
 \lim_{t\downarrow0}t E(t)=-\alpha M .
 \label{eq:E-limit}
\end{equation}

For the estimates, use \(\mathcal D=t\partial_t\) on profiles at fixed
\(y=x/L(t)\). Thus \(t\partial_t\) at fixed \(x\) is represented by
\(\mathcal D-\alpha y\partial_y\). Set \(\epsilon=1/64\) and
\(a=1+L^{-1}\phi\), and write \(\sigma=k-P\) in scaled coordinates.
Derivatives of \(\phi\) of any fixed order are bounded by
\(C_\epsilon\phi^{1-\epsilon}\). This follows by differentiating the
exponential near its support edge. For \(\Phi\) the same bounds
follow from \(\Phi'=-y\phi\) and \(\Phi\le\phi/2\).
Since \(\mathcal D\) multiplies each power of \(t\) by a constant,
\[
 |\mathcal D^j\sigma|\lesssim a,\qquad
 |\mathcal D^j\partial_y^h\sigma|
 \lesssim L^{-1}\phi^{1-\epsilon}\lesssim L^{-\epsilon}a
 \quad (j\le 2,\ 1\le h\le2).
\]
The latter derivatives also have the bound \(CL^{-1}\).
The identity \(g'=1/(1-mg^{m-1})\) gives \(|g'|\lesssim1\) and
\(|g''|\lesssim a^{-1}\). For \(J=g^m\), differentiation up to order
three yields \(|J^{(j)}(\sigma)|\lesssim a^{m-j}\).
Here \(g(\sigma)\asymp a\), and the second and third derivatives of
\(g\) are bounded by \(Cg^{m-2}\) and \(Cg^{m-3}\), as follows by
differentiating \(g'\).
We obtain on \(0\le x\le L\)
\begin{equation}
\begin{aligned}
 |U_t|+|\mathcal D U_t|&\lesssim (tL)^{-1},\\
 |(U^m)_{xx}|+|\mathcal D (U^m)_{xx}|&\lesssim L^{-2-m-2\epsilon},\\
 |\partial_x U^{m/2}|&\lesssim L^{-1-m/2-\epsilon}.
\end{aligned}
\label{eq:corebounds}
\end{equation}
In the first estimate use the bound \(CL^{-1}\) for \(\sigma_y\)
and \(\mathcal D\sigma_y\) in
\(U_t=t^{-1}g'(\sigma)(\mathcal D-\alpha y\partial_y)\sigma\).
In the second estimate the terms are
\(L^{-2}J''(\sigma)\sigma_y^2\) and
\(L^{-2}J'(\sigma)\sigma_{yy}\). The \(\mathcal D\)-derivative of
either term has the same bound because
\(\mathcal D\sigma=O(a)\). The last estimate uses
\((g^{m/2})'=O(a^{m/2-1})\).

Since \(\alpha(1+m+2\epsilon)<1\), \eqref{eq:corebounds} gives
\(|F|+|\mathcal D F|\lesssim (tL)^{-1}\). On the support,
\(\ell\asymp\log(1/t)\). Also
\begin{equation}
 |\log W-\ell|\lesssim t^\kappa,\qquad
 |\mathcal D\log U|+|\mathcal D\log W|\lesssim1
 \label{eq:logW}
\end{equation}
for some \(\kappa>0\). To see the first estimate, use \eqref{eq:heatbounds}
to compare \(B(t,Ly)\) with \(B(t,0)\). Their relative difference is
\(O(L/\sqrt t)\), and the relative size of \(1+P-U^m\) is
\(O(t^{\alpha+\beta-1}+t^{\beta-\alpha m})\).
For the second estimate,
\(\mathcal D B(t,Ly)=tB_t+\alpha x B_x=O(t^{-\beta})\).
The other terms in \(W\) and their scaled time derivatives are
\(O(1+\dot L+L^{-m})\). For \(U\), use \(\mathcal D\sigma=O(a)\)
and \(U\asymp a\).
It follows that \(|e|+|\mathcal D e|\lesssim\ell/(tL)\) on \(0\le x\le L\).

We now prove \eqref{eq:E-limit}. The diffusion part of \(E\) has
absolute value \(O(L^{-1-m-2\epsilon}\ell)=o(t^{-1})\).
Write \(U=1+k+R\), with \(R=U^m-1-P\) supported on \(0\le x\le L\).
The bounds on the scaled derivatives give
\[
 |R_t|\lesssim t^{-1}(L^{-m-\epsilon}+\dot L),\qquad
 \int_0^1 |R_t|\,dx
 \lesssim t^{-1}(L^{1-m-\epsilon}+L\dot L)=o((t\ell)^{-1}).
\]
Thus we may replace \(U_t\) by \(k_t\) against either logarithm in
\(E\), at an error \(o(t^{-1})\). The term \(\int k_t\log W\) is
\(o(t^{-1})\) by \(\int k_t=0\), \eqref{eq:logW}, and
\(\int |k_t|=O(t^{-1})\). In the other term, \(\log U-\log(1+k)\)
is uniformly bounded and tends to zero for each fixed \(y\) with
\(\phi(y)>0\). These facts follow from \(U\asymp1+k\),
\(m<1\), and \(P/k=O(L\dot L)\) where \(k>0\).
Since \(t k_t\,dx=-\alpha(y\phi)'dy\), dominated convergence
allows replacement by \(\log(1+k)\). Integration by parts gives
\[
\begin{aligned}
 t\int_0^1 k_t\log(1+k)\,dx
 &=-\alpha\int_0^1(\phi+y\phi')\log(1+\phi/L)\,dy\\
 &=\alpha\int_0^1 y\phi'\frac{\phi}{L+\phi}\,dy
 \longrightarrow\alpha\int_0^1y\phi'\,dy=-\alpha M .
\end{aligned}
\]
The boundary term is zero at \(y=0,1\). This proves \eqref{eq:E-limit}.

\subsection{A third concentration and the pointwise entropy bound}

To use the negative integral in \eqref{eq:E-limit}, we introduce a
third concentration whose reaction term redistributes \(e\) across the
support. Let \(\rho\) be an even smooth function, supported on
\([-2,2]\), positive on \((-2,2)\), nonincreasing for \(y\ge0\), and
normalized by \(\int_0^2\rho(y)\,dy=1\). Set
\(\rho_L(x)=L^{-1}\rho(x/L)\), and restrict \(t\) so that
\(2L<1\). Define
\begin{equation}
 S=-\frac{e}{\ell}+\frac{E}{\ell}\rho_L,\qquad
 K(t,x)=-\int_x^1(s-x)S(t,s)\,ds,\qquad V=B+K .
 \label{eq:reservoir}
\end{equation}
The choice of \(S\) gives \(\int S=0\), so \(-K_{xx}=S\) with
\(K_x(0)=K_x(1)=0\). Also \(K=0\) for \(x\ge2L\).
Evenness of \(S\) at zero and \(K_x(0)=0\) give a smooth even
extension of \(K\). The reaction residual of \(V\) is \(G=K_t+S\).

Writing \(E=L\int_0^1 e(t,Ly)\,dy\) gives
\(|E|+|tE'|\lesssim \ell/t\), and \(t|\ell'|\lesssim1\) by
\eqref{eq:heatbounds}. Since \(\rho\) has a positive minimum on
\([0,1]\),
\[
 |S(t,Ly)|+|\mathcal D S(t,Ly)|\lesssim (tL)^{-1}\rho(y)
 \qquad (0\le y\le 2).
\]
In scaled coordinates \(K=-L^2\int_y^2(s-y)S(t,Ls)\,ds\).
Using monotonicity of \(\rho\) in this integral and in its \(y\)
derivative gives
\begin{equation}
 |K|+|\mathcal D K|\lesssim (L/t)\rho(y),\qquad
 |K_x|\lesssim t^{-1}\rho(y),\qquad
 |K_t|\lesssim (L/t^2)\rho(y).
 \label{eq:Kbounds}
\end{equation}
For the last bound apply
\(t\partial_t|_x=\mathcal D-\alpha y\partial_y\).
In particular these estimates retain the factor \(\rho(y)\) even
at the edge of its support.

Since \(\alpha+\beta>1\) and \(2L\le\sqrt t\) for small \(t\),
\(|K|/B=O(t^{\alpha+\beta-1})\) on \(0\le x\le 2L\). Hence
\(V\ge B/2\), and \(V\ge c_0t^{-\beta}\) there for some \(c_0>0\).
We also have \(|\log V-\ell|=O(t^\kappa)\) on this support for
some \(\kappa>0\), and \(V\ge2\) there for small \(t\).
The three reaction residuals now obey the pointwise inequality
\begin{equation}
\begin{aligned}
 e+G\log V
 &=E\rho_L+S(\log V-\ell)+K_t\log V\\
 &\le \frac{\rho(y)}{tL}
   \bigl(-c+C t^\kappa+C(L^2/t)\ell\bigr)\le0
 \qquad (0\le x\le2L).
\end{aligned}
\label{eq:pointwiseentropy}
\end{equation}
Here \(E\le-c/t\) by \eqref{eq:E-limit}, and
\(L^2\ell/t\to0\) because \(2\alpha>1\).
The residuals are zero for \(x\ge2L\).

\subsection{Energy bounds for the diffusion powers}

The identity \(U+W=1+B+k\) and positivity give uniform \(L^1\)
bounds on \(U,W\). From \eqref{eq:Kbounds},
\(\int|K|\,dx\lesssim L^2/t\), so the same holds for \(V\).
The spatial integral of \(|\partial_x U^{m/2}|^2\) is
\(O(L^{-1-m-2\epsilon})\) by \eqref{eq:corebounds}, an integrable
function of time. For the other two components, use
\(|P_x|=(\dot L/L)|y\phi(y)|\lesssim t^{-1}\) from \eqref{eq:spike},
the bound on \(K_x\) in \eqref{eq:Kbounds}, and the lower bound on \(B\).
These give
\[
 \int_0^{2L}\frac{|P_x|^2+|K_x|^2}{B}\,dx
 \lesssim L t^{\beta-2}
 \quad\text{and}\quad
 \frac{|(U^m)_x|^2}{B}
   =4\,\frac{U^m}{B}\,|\partial_x U^{m/2}|^2 .
\]
The first right-hand side is time-integrable by
\(\alpha+\beta>1\), and \(U^m/B\) is bounded on the support
of its derivative. Combining these estimates with
\eqref{eq:heatenergy} and \(W,V\ge B/2\) bounds the two square-root
gradient energies. The \(L^2\) norms of the half-powers follow
from the mass bounds and \(U^m\le1+U\).

\section{The reaction and the initial trace}

\subsection{An autonomous reaction}

Choose \(t_*>0\) small enough that all preceding estimates and inequalities
hold for \(0<t\le2t_*\). To realize the residuals as a function of the
concentrations, take two further linear components
\[
 A(t,x)=1+t,\qquad X(x)=2+\cos(\pi x).
\]
Let \(\eta_4,\eta_5\) be smooth cutoffs with compact supports in
\((0,\infty)\) and equal to one on \([1,1+t_*]\) and \([1,3]\),
respectively. Define \(f_4(z)=\eta_4(z_4)\) and
\(f_5(z)=\pi^2(z_5-2)\eta_5(z_5)\). These give the residuals of
\(A,X\) on \(0<t\le t_*\).

For the first three components of \(f(z)\), put \(\tau=z_4-1\).
Use \(x(z_5)=X^{-1}(z_5)\) on \([1,3]\), continued constantly
as \(1\) on \(z_5\le1\) and \(0\) on \(z_5\ge3\).
When \(0<\tau<2t_*\), evaluate the profiles \(F,G\) at
\((\tau,x(z_5))\). Choose smooth cutoffs valued in \([0,1]\)
with the following properties. The cutoff \(\theta\) equals one on
\([0,t_*]\) and vanishes near \([2t_*,\infty)\). Each
\(\zeta_1,\zeta_2\) is zero below and at \(1/2\), and equals one
at arguments at least \(1\). The cutoff \(\zeta_3\) is zero
below and at \(3/2\), and equals one at arguments at least \(2\).
Finally, \(\chi\) is zero below and at \(c_0/2\), and equals one
at arguments at least \(c_0\), with \(c_0\) from the lower bound on \(V\).
Set
\[
 \lambda=\theta(\tau)\zeta_1(z_1)\zeta_2(z_2)
           \zeta_3(z_3)\chi(\tau^\beta z_3).
\]
Where \(0<\tau<2t_*\) and \(\lambda>0\), define
\begin{equation}
\begin{aligned}
 H_z&=F(\log z_1-\log z_2)+G\log z_3,\\
 (f_1,f_2,f_3)&=\lambda
       \left(F,-F,G-\frac{(H_z)_+}{\log z_3}\right),
\end{aligned}
\label{eq:reaction}
\end{equation}
where \(r_+=\max(r,0)\). Elsewhere set \((f_1,f_2,f_3)=0\).

This is a locally Lipschitz function on the whole closed cone.
On a compact positive-time strip in \(0<\tau<2t_*\), \(F\) and \(G\)
are smooth and even at the spatial endpoints. Their \(x\)-derivatives
are therefore bounded by \(C\operatorname{dist}(x,\{0,1\})\).
Since \(|X'(x)|=\pi\sin(\pi x)\) is bounded below by a positive
multiple of this distance, the compositions with the clamped inverse
are Lipschitz in \(z_5\), uniformly on the strip. Their time
derivatives are bounded there, which gives joint local Lipschitz
continuity. On such a strip one can replace \(z_1,z_2,z_3\) inside
the logarithms of \eqref{eq:reaction} by their maxima with
\(1/2,1/2,3/2\), respectively. These values coincide with the
original arguments where \(\lambda>0\) and keep the denominator
positive. The resulting locally Lipschitz product also gives the
specified zero value where \(\lambda=0\). At a finite state with
\(\tau=0\), \(z_3\)
is locally bounded, and \(\chi(\tau^\beta z_3)\) vanishes in a
right neighborhood of zero. This gives local Lipschitz continuity
across \(\tau=0\). The cutoff \(\theta\) gives the corresponding
zero neighborhood at \(2t_*\).

The reaction vanishes in its own component on each zero-coordinate
face, hence is quasi-positive. The relation \(f_1+f_2=0\) holds on the
whole cone. On the positive cone the first three terms in
\(\sum_i f_i\log z_i\) sum to \(\lambda(H_z-(H_z)_+)\le0\) on
the active region and to zero outside. The last two terms are
bounded, and \(1+s\log s\ge1-\mathrm e^{-1}\) for \(s>0\). Thus
\eqref{eq:entropy} holds with all \(\mu_i=0\) and a fixed finite
\(C\ge0\).

Along \(u=(U,W,V,A,X)\) for \(0<t\le t_*\), the time and space recovered
from \(z_4,z_5\) are the actual \(t,x\). Where \(F,G\) may be nonzero,
we have \(U\ge1\), \(W,V\ge2\), and \(t^\beta V\ge c_0\), so
\(\lambda=1\). The inequality \eqref{eq:pointwiseentropy} gives
\(H_z\le0\) along the profiles. Therefore the first three components
of the reaction are exactly \(F,-F,G\) there. Outside their support
both \(F\) and \(G\) vanish. All five components solve
\eqref{eq:system} classically on each positive-time strip up to \(t_*\)
with zero normal derivative of their diffusion potentials.

\subsection{The truncated initial terms}

Take the assigned initial data
\begin{equation}
 u_0=(1,b_0,b_0,1,2+\cos(\pi x)).
 \label{eq:initial}
\end{equation}
Every component has finite initial entropy. At each fixed
\(x\in(0,1)\), eventually \(x>2L(t)\), and the five fields there are
\((1,B,B,1+t,X)\). This proves \(u(t,x)\to u_0(x)\) almost everywhere.
The first three components have the required mass and energy bounds
proved above. The last two are smooth and bounded up to time zero.

Fix a permitted \(\xi,\psi\) and \(0<\varepsilon<T\le t_*\).
On \((\varepsilon,T)\) the classical chain rule and spatial
integration by parts give \eqref{eq:renormalized} with initial term
at \(\varepsilon\). Write \(w_i=u_i^{m_i/2}\) as in
\eqref{eq:renormalized}. The concentrations are positive on this strip, so
\[
 u_i^{m_i-1}\nabla u_i=\frac{2}{m_i}w_i\nabla w_i,\qquad
 \nabla u_j=\frac{2}{m_j}u_j^{1-m_j/2}\nabla w_j .
\]
These yield exactly the coefficients in \eqref{eq:renormalized}.
The no-flux conditions eliminate all spatial boundary terms.

The function \(\xi\) is bounded, as it is constant along each ray
after that ray leaves a ball containing \(\operatorname{supp} D\xi\).
Thus bounded convergence gives the initial term with \eqref{eq:initial}
as \(\varepsilon\downarrow0\). The terms involving \(D\xi,D^2\xi\)
have bounded concentration weights on their compact supports.
The gradient product is integrable by the two \(L^2\) gradient
bounds, and the single-gradient flux term is integrable on the finite
cylinder. The reaction term is integrable since \(f\) is bounded on the
support of \(D\xi\). The time test term converges by boundedness of
\(\xi\). This proves \eqref{eq:renormalized} for every
\(T\in(0,t_*]\).

\subsection{Continuation and the mass difference}

Theorem 1.3 of Fellner, Fischer, Kniely and Tang \cite{FFKT} gives a
global renormalized solution starting from \(u(t_*)\). Indeed
\((0,1)\) is bounded Lipschitz, all diffusion exponents are in
\((0,2)\), and \(u(t_*)\) is bounded and positive with finite entropy.
Their conditions (F1)--(F3) are local Lipschitz continuity,
quasi-positivity, and the logarithmic reaction bound \eqref{eq:entropy}.
The bound extends to zero concentrations by giving \(f_i\log z_i\) value
zero at \(z_i=0\), because each \(f_i\) vanishes in a neighborhood of
that face. Their Definition 1.1 gives the truncated identity (1.4) and
the \(L^\infty_tL^1_x\) bound on every finite cylinder. Remark 1.4
gives the \(L^2_tH^1_x\) bound for the diffusion half-powers of the
constructed solution. Write \(v\) for this continuation, with time
starting at zero.

Keep \(u\) on \((0,t_*]\), and take \(v(t-t_*)\) for \(t>t_*\).
For a smooth \(\psi\) and almost every \(T>t_*\), apply the
continuation identity on \((0,T-t_*)\) with the translated test.
Add it to the identity for the first segment at \(t_*\).
The two terms \(\int_0^1 \xi(u(t_*))\psi(t_*)\,dx\) cancel by the
prescribed initial datum of \(v\). This gives
\eqref{eq:renormalized} on the whole cylinder, and the two segments
give its regularity requirements for every finite \(T\).

On \(0<t\le t_*\), the first two concentrations satisfy
\(U+W=B+1+k\). Therefore
\[
 \int_0^1(U+W)\,dx
 =1+\int_0^1 b_0\,dx+M,\qquad
 \int_0^1(u_{1,0}+u_{2,0})\,dx=1+\int_0^1 b_0\,dx .
\]
Since \(M>0\), this proves Theorem~\ref{thm:counterexample}.

\section*{AI-use disclosure}
OpenAI Codex assisted in developing the construction, checking the argument,
and drafting and reviewing the manuscript.

\begin{thebibliography}{9}

\bibitem{AIM483}
AIM problem collection.
\newblock Problem 483, \emph{Automatic conservation laws for renormalized
nonlinear diffusion}.
\newblock \href{https://github.com/MColbrook/AIM/blob/b2ec0fac559f532fc0ed04fce81deea0ce07c68b/problems/483-nonlinear-diffusion-renormalization-mass-conservation.md}{Pinned problem statement}.

\bibitem{FFKT}
K.~Fellner, J.~Fischer, M.~Kniely, and B.~Q.~Tang.
\newblock Global renormalised solutions and equilibration of
reaction--diffusion systems with nonlinear diffusion.
\newblock \emph{Journal of Nonlinear Science}, 33, article~66, 2023.
\newblock \href{https://doi.org/10.1007/s00332-023-09926-w}{doi:10.1007/s00332-023-09926-w}.

\bibitem{Fischer}
J.~Fischer.
\newblock Weak--strong uniqueness of solutions to entropy-dissipating
reaction--diffusion equations.
\newblock \emph{Nonlinear Analysis}, 159:181--207, 2017.
\newblock \href{https://doi.org/10.1016/j.na.2017.03.001}{doi:10.1016/j.na.2017.03.001}.
\newblock Author preprint \href{https://arxiv.org/abs/1703.00730v1}{arXiv:1703.00730v1}.

\end{thebibliography}
\end{document}
